% =============================================================================
% Proposed Solution - Execute (EX) stage
% Author: Kapish Rohilla (2401010210)
% =============================================================================
\subsection{Execute Stage}
\label{sec:ex_stage}

The Execute (EX) stage performs all arithmetic and logic operations, computes
load/store addresses, and resolves branches and jumps. It is implemented by the
wrapper module \texttt{ex\_stage}, which contains the operand forwarding
multiplexers and instantiates three sub-modules: \texttt{alu\_control},
\texttt{alu} and \texttt{branch\_unit}. Its results are latched into the
EX/MEM pipeline register (\texttt{ex\_mem\_reg}). Instruction encodings follow
the RISC-V unprivileged specification~\cite{kapish_riscv_isa}, and the overall
organisation follows the classic five-stage pipeline described
in~\cite{kapish_patterson_riscv, kapish_harris_riscv}.

\subsubsection{Arithmetic Logic Unit}
The ALU is purely combinational. A 4-bit control code selects one of ten
operations (Table~\ref{tab:alu_ops}). Shift amounts use only the low five bits
of operand $B$, as required for RV32I. The \texttt{zero} output is asserted
when the result is zero; branch decisions, however, are taken in the branch
unit, not from this flag.

\begin{table}[htbp]
\caption{ALU operations}
\label{tab:alu_ops}
\centering
\begin{tabular}{|c|l|l|}
\hline
\textbf{Code} & \textbf{Operation} & \textbf{Result} \\
\hline
0 & ADD  & $A + B$ \\
1 & SUB  & $A - B$ \\
2 & AND  & $A \,\&\, B$ \\
3 & OR   & $A \mid B$ \\
4 & XOR  & $A \oplus B$ \\
5 & SLL  & $A \ll B[4{:}0]$ \\
6 & SRL  & $A \gg B[4{:}0]$ (zero fill) \\
7 & SRA  & $A \gg B[4{:}0]$ (sign fill) \\
8 & SLT  & $1$ if $A < B$ (signed), else $0$ \\
9 & SLTU & $1$ if $A < B$ (unsigned), else $0$ \\
\hline
\end{tabular}
\end{table}

\subsubsection{ALU Control}
The main control unit in ID produces a 2-bit \texttt{alu\_op}. The ALU control
unit combines it with \texttt{funct3} (instr[14:12]) and \texttt{funct7[5]}
(instr[30]) to produce the 4-bit ALU code (Table~\ref{tab:alu_ctrl}).
For I-type instructions, bit 30 belongs to the immediate, so \texttt{funct3}
$=000$ is always decoded as ADDI and never as SUB; bit 30 is used only to
distinguish SRAI from SRLI.

\begin{table}[htbp]
\caption{ALU control decoding}
\label{tab:alu_ctrl}
\centering
\begin{tabular}{|c|c|c|l|}
\hline
\textbf{alu\_op} & \textbf{funct3} & \textbf{funct7[5]} & \textbf{ALU operation} \\
\hline
00 & x   & x & ADD (LW, SW, LUI) \\
01 & x   & x & SUB (branches) \\
\hline
10 & 000 & 0/1 & ADD / SUB \\
10 & 001 & x & SLL \\
10 & 010 & x & SLT \\
10 & 011 & x & SLTU \\
10 & 100 & x & XOR \\
10 & 101 & 0/1 & SRL / SRA \\
10 & 110 & x & OR \\
10 & 111 & x & AND \\
\hline
11 & 000 & x & ADD (ADDI) \\
11 & 001 & x & SLL (SLLI) \\
11 & 010 & x & SLT (SLTI) \\
11 & 011 & x & SLTU (SLTIU) \\
11 & 100 & x & XOR (XORI) \\
11 & 101 & 0/1 & SRL / SRA (SRLI / SRAI) \\
11 & 110 & x & OR (ORI) \\
11 & 111 & x & AND (ANDI) \\
\hline
\end{tabular}
\end{table}

\subsubsection{Operand Forwarding Multiplexers}
To avoid stalling on read-after-write dependences, each register operand passes
through a 3-input multiplexer controlled by the forwarding unit:
\texttt{2'b10} selects the EX/MEM result, \texttt{2'b01} selects the
write-back value from MEM/WB, and \texttt{2'b00} selects the value read from
the register file in ID. EX/MEM has priority because it holds the most recent
value. After forwarding, operand $A$ is forced to zero for LUI (so the ALU
computes $0 + \mathit{imm}$), and operand $B$ is replaced by the immediate when
\texttt{alu\_src} is set. The forwarded $rs2$ value is also passed on as
\texttt{store\_data}, so a store always writes the latest value of its source
register.

\subsubsection{Branch and Jump Resolution}
The branch unit evaluates BEQ, BNE, BLT and BGE on the forwarded operands
(BLT and BGE use signed comparison). It asserts
\[
\mathit{pc\_src} = (\mathit{branch} \wedge \mathit{cond}) \vee \mathit{jump} \vee \mathit{jalr}
\]
and computes the target as $\mathit{pc} + \mathit{imm}$ for branches and JAL,
or $(\mathit{rs1} + \mathit{imm})$ with bit~0 cleared for JALR, as required
by the specification~\cite{kapish_riscv_isa}. The pipeline predicts every
branch as not taken. Because the decision is made in EX, two younger
instructions have already entered IF/ID and ID/EX when a branch or jump is
taken; the hazard unit flushes both registers, giving a two-cycle penalty.
Resolving in EX keeps the ID stage simple and lets branches use forwarded
operands without extra comparison hardware in ID~\cite{kapish_patterson_riscv}.

\subsubsection{Link Address for JAL and JALR}
JAL and JALR must write the return address $\mathit{pc}+4$ into $rd$. Instead
of adding a separate write-back path, the EX stage selects
\[
\mathit{ex\_result} = (\mathit{jump} \vee \mathit{jalr})\ ?\ \mathit{pc}+4 : \mathit{alu\_result}.
\]
The link address therefore travels through EX/MEM and MEM/WB exactly like an
ALU result, uses the existing write-back multiplexer, and can be forwarded to
a dependent instruction by the normal forwarding paths.

\subsubsection{EX/MEM Pipeline Register}
The \texttt{ex\_mem\_reg} module latches \texttt{reg\_write},
\texttt{mem\_read}, \texttt{mem\_write}, \texttt{mem\_to\_reg},
\texttt{ex\_result}, \texttt{store\_data} and \texttt{rd} on each rising clock
edge and clears them on synchronous reset. It needs no flush input: stalls and
flushes are handled by inserting a bubble into ID/EX, which then flows through
EX/MEM as an instruction with all control signals zero.

\subsubsection{Verification}
Each module has a self-checking Icarus Verilog testbench that prints PASS/FAIL
per check. The tests cover every ALU operation (including negative operands,
SLT versus SLTU, SRA sign extension and shifts by 31), the ALU-control code
for every supported instruction, each branch taken and not taken (including
signed negative operands), JAL and JALR targets, all forwarding selections on
both operands, the LUI path, the JAL/JALR link value and the EX/MEM register
reset and pass-through behaviour.
